\subsection{The collage test}

\begin{proposition}[Proved: Lead 1, including the unrestricted homotopy type]\label{B:collage}
Let $A$ be small and $K:A\to\Set$ covariant. Form the category $A_K$ by adjoining an object $x$, with $A_K(x,a)=K(a)$, $A_K(a,x)=\varnothing$, and $A_K(x,x)=\{\id_x\}$. Composition of $h:x\to a$ with $v:a\to b$ is $K(v)h$. Put $E=\El_A(K)$ and let $\pi:E\to A$ forget the element. Then $i:A\hookrightarrow A_K$ is final if and only if $E$ is nonempty and connected. Without any contractibility assumption on $BA$, there is a natural homotopy-pushout description
\[
 B(A_K)\simeq\operatorname{hocolim}\bigl(BA\xleftarrow{B\pi}BE\longrightarrow *\bigr).
\]
More precisely, the displayed realization is the mapping cone of the specified map $B\pi:BE\to BA$, not a construction determined only by the two abstract homotopy types.
\end{proposition}

\begin{proof}
For $a\in A$, the category $a\backslash A$ for the inclusion has initial object $(a,\id_a)$. For the new object $x$, the comma category is precisely $\El_A(K)$. This proves finality with the stated edge cases.

Write $CE=\Delta^0*NE$ for the simplicial cone with its new vertex first. There is an isomorphism of simplicial sets
\[
 N(A_K)\cong NA\amalg_{NE}(\Delta^0*NE).
\]
To check it, a nondegenerate simplex outside $NA$ starts with $x$, followed by an arrow $h:x\to a_0$ and a string $a_0\to\cdots\to a_n$. The data are exactly a simplex of $NE$, because the remaining elements of the string are determined from $h$ by $K$. Deleting the first vertex forgets $h$ and is the map $N\pi$; all other faces and degeneracies agree with the cone. This also proves the assertion for repeated identity arrows. Realization preserves this pushout, and $BE\hookrightarrow CBE$ is a CW cofibration. Replacing the contractible cone by a point identifies it with the indicated homotopy pushout. If $E$ is empty, the result is $BA\amalg *$, as it should be.
\end{proof}

\begin{corollary}[Proved: the integral obstruction and its converse for one collage]\label{B:acyclic-test}
Suppose $BA$ is weakly contractible and $E=\El_A(K)$ is nonempty and connected. Then
\[
 B(A_K)\simeq\Sigma BE,
 \qquad
 \widetilde H_n(B(A_K);\mathbb Z)\cong
 \widetilde H_{n-1}(BE;\mathbb Z)
 \quad(n\geq1).
\]
Consequently $B(A_K)$ is weakly contractible if and only if $BE$ is integrally acyclic. In particular, if $T_{\Winfty}(A)$, every such $E$ is acyclic. The same necessary assertion with coefficients in a commutative ring $R$ holds under $T_{W_R}(A)$.
\end{corollary}

\begin{proof}
The map $BE\to BA$ has contractible target, so its homotopy cofiber is the unreduced suspension of $BE$. The suspension of a connected CW complex is simply connected, and its reduced homology is the shifted reduced homology of that complex. A simply connected acyclic CW complex is weakly contractible by the Hurewicz theorem, applied at the least hypothetical nonzero homotopy group. For $T_{\Winfty}(A)$, both $A$ and its final image $A_K$ are weakly contractible. For $W_R$, use the homology long exact sequence of the same cofiber; no assertion of simple connectivity is needed in that case.
\end{proof}

\begin{proposition}[Proved: finite and transfinite iteration]\label{B:iteration}
If $T_W(A)$ and $\El_A(K)$ is nonempty and connected, then $T_W(A_K)$. This remains true after any small transfinite succession of such one-object collages, taking unions at limit ordinals. Thus every connected covariant element-category test on every category obtained in this way is a necessary test, not just those on the starting category.
\end{proposition}

\begin{proof}
The successor assertion is Corollary~\ref{A:final-images}. For completeness, the inclusion of an earlier stage into a union at a limit ordinal is final. If $b$ first appears at a stage $\beta$, its comma category into a fixed earlier stage becomes nonempty connected by induction on $\beta$: successor steps are composition of final inclusions, and at a limit the comma category is the increasing union of the earlier comma categories. Once nonempty, all are connected, so their union is connected. Equivalently, for the inclusion of the original stage, an object $b$ belongs to some stage $\beta$, and no later collage adds arrows from an existing object to an old object, so $b\backslash A$ is already the comma category at that stage. The original inclusion into the full union is therefore final, and Corollary~\ref{A:final-images} applies. The smallness assumption ensures that the union is still a category in our universe.
\end{proof}

The element-category characterization in Section C gives a closed form for these iterated necessary conditions in the integral and weak-homotopy cases. It avoids any inference from isomorphisms of ind-objects or from unspecified weak equivalences between stages.

\subsection{The parallel-pair homology formula}

\begin{proposition}[Proved: the integral formula in Lead 2]\label{B:lead2}
Let $a,b:p\to q$ be a parallel pair in any small category $A$. Let $h_p=A(p,-)$ and let $H$ be the pointwise quotient of $h_p$ generated by $ca\sim cb$, for every $c:q\to r$ and every $r\in A$. Then $\operatorname{El}_A(H)$ is nonempty and connected and
\[
 H_1(B\operatorname{El}_A(H);\mathbb Z)\cong\mathbb Z/I,
\]
where $I\subseteq\mathbb Z$ is the subgroup generated by all integers $\sum_cn_c$ such that
\[
 \sum_cn_c(ca-cb)=0\quad\text{in }\mathbb Z[A(p,r)]
\]
for one object $r$ and a finitely supported family $(n_c)_{c\in A(q,r)}$. Equivalently, if $I=g\mathbb Z$ with $g\geq0$, this group is $\mathbb Z/g\mathbb Z$. This formula has no asphericity assumption on $A$.
\end{proposition}
\begin{proof}
Put $P_s=\mathbb Z[A(s,-)]\in\Ab^A$. These are projective: $\operatorname{Hom}_{\Ab^A}(P_s,F)=F(s)$ naturally, and evaluation is exact. Define $d:P_q\to P_p$ by $d(c)=ca-cb$, let $K=\ker d$ and $J=\operatorname{im}d$. The free-abelian-group functor preserves coequalizers, so we have exact sequences
\[
 0\to J\to P_p\to\mathbb Z[H]\to0,
 \qquad 0\to K\to P_q\to J\to0.
\]
Both $\colim P_q$ and $\colim P_p$ identify canonically with $\mathbb Z$, and $\colim d$ is zero. The long exact sequence of derived colimits for the first sequence, and right exactness applied to the second, give
\[
 L_1\colim_A\mathbb Z[H]
 \cong\colim_A J
 \cong\operatorname{coker}(\colim_AK\to\mathbb Z).
\]
For the first identification, the map $\colim J\to\colim P_p$ is zero because it is zero after the epimorphism $\colim P_q\to\colim J$. The image of $\colim K\to\mathbb Z$ is exactly the subgroup $I$ displayed in the statement.

The bar chain complex for $\mathbb Z[H]$ has basis a string $a_0\to\cdots\to a_j$ and an element of $H(a_0)$. These are exactly the simplices of $N\operatorname{El}_A(H)$, compatibly with every face. It therefore computes the homology on the left. Finally, every object of $\operatorname{El}_A(H)$ receives an arrow from $(p,[\id_p])$: choose a representative $f:p\to r$ of its element. Thus this category is connected and nonempty, although the arrow from that object need not be unique.
\end{proof}

\subsection{The smallest band obstruction}

\begin{proposition}[Proved: Lead 3]\label{B:right-zero}
Let $M=\{1,e,f\}$, with $1$ a unit and $uv=v$ for $u,v\in\{e,f\}$. Regard $M$ as a one-object category with multiplication $uv=u\circ v$. It is cofiltered, its classifying space is contractible, and it is not filtered. Let $K:M\to\Set$ have underlying set $\{h,i\}$, with $1$ acting identically and both $e,f$ acting constantly with value $i$. Then
\[
 B\El_M(K)\simeq S^1,
 \qquad B(M_K)\simeq S^2.
\]
The inclusion $M\to M_K$ is final. In particular $M$ is not totally $\Winfty$-aspherical and is not totally $W_R$-aspherical for any nonzero commutative ring $R$.
\end{proposition}

\begin{proof}
Associativity follows from the right-zero rule, with the unit inserted. For any $a,b\in M$, $ae=be=e$, proving cofilteredness. No $c\in M$ satisfies $ce=cf$, since $ce=e$ and $cf=f$. This disproves filteredness. The element $e$ defines a natural transformation from the constant endofunctor of the one-object category (sending every arrow to $1$) to its identity: the naturality equation is $me=e$. Hence $BM$ is contractible.

The action is a functor because compositions of the two constant functions are that same constant function. Its element category has two objects, $h$ and $i$. There are only the identity at $h$, two arrows $h\rightrightarrows i$ labelled $e,f$, and the three endomorphisms $1,e,f$ at $i$. Postcomposition fixes each of the two arrows $h\to i$, by the right-zero multiplication. Therefore this element category is itself the collage of the covariant constant two-element functor on the full one-object category at $i$. The category of elements of that constant functor is the disjoint union of two copies of $M$. Proposition~\ref{B:collage} now computes its classifying space as the homotopy cofiber of
\[
 BM\amalg BM\longrightarrow BM,
\]
which is $S^1$, since each copy of $BM$ is contractible. It is connected, so $M\to M_K$ is final. A second application of the proposition gives the suspension $S^2$. Its degree-two homology with nonzero coefficients $R$ is $R$.
\end{proof}

The complete multiplication and action tables, all arrows and composites in the two collages, and normalized nerve-chain computations through degree four are recorded in the accompanying program \path{code/root_monoid_checks.py} and its output. The all-degree assertion above follows from the two homotopy-pushout calculations; a truncated chain calculation alone would not prove it.

\begin{proposition}[Proved: the whole rectangular-band family]\label{B:bands}
Let $I,J$ be nonempty small sets, let $R_{I,J}=I\times J$ with $(i,j)(k,l)=(i,l)$, and let $M_{I,J}=R_{I,J}^{1}$ have a new identity adjoined. Then $M_{I,J}$ is filtered, and totally $\Winfty$-aspherical, if and only if $|J|=1$. It is cofiltered if and only if $|I|=1$. For every nonzero $R$, total $W_R$-asphericity has the same classification. If the band is empty, adjoining an identity gives the trivial monoid and all these properties hold.
\end{proposition}

\begin{proof}
If $J$ is a singleton, any band element is a left zero and postcomposition by it equalizes all arrows; Theorem~\ref{A:filtered} applies. If $J$ has at least two elements, choose a surjection $J\twoheadrightarrow\{e,f\}$ and forget the $I$ coordinate. This gives a surjective homomorphism $M_{I,J}\to M$ to Proposition~\ref{B:right-zero}'s monoid. Any homomorphism surjective on the underlying set is a full identity-on-objects functor. Its unique comma category is connected, since its identity object has an arrow to every object via a lift of the corresponding target monoid element. Hence it is final, and totality would pass to $M$, a contradiction. This works equally for all nonzero $R$. If $|J|>1$, two band elements with different second coordinates cannot be equalized by left multiplication; this also checks filteredness directly. The cofiltered assertion is the dual argument on first coordinates. No finiteness assumption is used.
\end{proof}

\subsection{Other final tests}

\begin{proposition}[Proved: quotient, localization, and idempotent-completion tests]\label{B:other-final}
Every quotient $q:A\to A/{\sim}$ by a congruence on the hom-sets, with unchanged objects, is final. Every localization $A\to A[S^{-1}]$ at a set of arrows is final. The canonical fully faithful functor $j:A\to\Kar(A)$ into its idempotent completion is final. Consequently a totally $W$-aspherical category has totally $W$-aspherical categories as all these specified images.
\end{proposition}

\begin{proof}
For a congruence quotient, an object of $b\backslash A$ is an arrow $b\to a$ in the quotient. It lifts to an arrow in $A$, giving a morphism from $(b,\id_b)$ to that object. The comma category is nonempty connected.

For a localization, an arrow $b\to a$ is represented by a finite word of original arrows and inverses of arrows in $S$. Start at $(b,\id_b)$ and read the word. A forward arrow gives a forward morphism in $b\backslash A$. An inverse $s^{-1}$ gives a backward morphism using the original arrow $s$ of $A$. This constructs a zigzag to $(a,\alpha)$, irrespective of different representatives of $\alpha$. Thus the comma category is connected and nonempty.

An object of $\Kar(A)$ is $(a,e)$ with $e^2=e$; an arrow $(a,e)\to(b,d)$ is $v:a\to b$ with $v=dve$. The comma category $(a,e)\backslash A$ has distinguished object $(a,e:(a,e)\to j(a))$. Any object $(b,v)$ receives a morphism from it, namely $v:a\to b$, since $ve=v$. This again proves connectedness. Applying Corollary~\ref{A:final-images} gives the consequences.
\end{proof}

\begin{proposition}[Proved: totality is invariant under idempotent completion]\label{B:karoubi}
For every fundamental localizer $W$, $T_W(A)$ holds if and only if $T_W(\Kar(A))$.
\end{proposition}

\begin{proof}
Only the converse needs proof. Suppose $\Kar(A)$ is total and $f:A\to B$ is final. We claim that $\Kar(f):\Kar(A)\to\Kar(B)$ is final. Fix $(b,d)\in\Kar(B)$. Every arrow $\alpha:(b,d)\to\Kar(f)(a,e)$ gives a comma morphism to the object $(j(a),\alpha)$ via the splitting arrow $(a,e)\to j(a)$. For objects of this last form, finality of $f$ connects the corresponding arrows in $b\backslash A$. Composing every arrow of that connecting zigzag on the right with $d$ keeps it in the subcategory of arrows $\alpha$ satisfying $\alpha d=\alpha$. Nonemptiness follows by choosing any $b\to fa$ and precomposing with $d$. Thus $\Kar(f)$ is final.

It follows that $\Kar(B)$ is $W$-aspherical. The functor $B\to\Kar(B)$ belongs to $\Winfty$: for its comma category considered above, the distinguished object maps naturally to every object, with components the underlying $v$, so there is a natural transformation from its constant functor to the identity. These comma categories are contractible, and the dual form of Quillen's Theorem A applies. Cisinski's minimality theorem places this functor in $W$. Two-out-of-three gives $B\to e\in W$, proving $T_W(A)$.
\end{proof}

\begin{remark}[Proved: what the tests actually force]\label{B:scope-tests}
A total category must therefore remain aspherical after every hom-set congruence and every localization, including its groupoid completion. Ordinary quotient tests do not require a faithful comparison map. Idempotent completion is stronger: it both preserves and reflects totality by the specified functors in Proposition~\ref{B:karoubi}. None of these assertions says that an arbitrary weak equivalence preserves totality. The common-successor obstruction, the affine parallel-arrow obstruction of Lead 2, and the complete element-category criterion appear in Section C.
\end{remark}

\begin{example}[Proved: the infinite fence and the nonuniformity of contractions]\label{B:fence}
Let $F$ have objects $\mathbb N$ and nonidentity arrows $2r\to2r+1$ for $r\geq0$ and $2r\to2r-1$ for $r\geq1$, with no other nonidentity arrows. Then $BF$ is contractible, but $F$ is not totally $\Winfty$-aspherical. Its identity endofunctor is not connected to a constant endofunctor by any finite zigzag of natural transformations.
\end{example}
\begin{proof}
There are no composable pairs of nonidentity arrows, and the realization of the nerve is exactly an infinite ray. The maximal objects $1$ and $3$ have no common successor, so Lemma~\ref{C:directed-affine} rules out totality. A component of any natural transformation between two functors into $F$ moves its object by at most one in the undirected path metric. A zigzag of length $m$ therefore moves each object by at most $m$. A constant endofunctor with value $c$ would require moving the object $n$ through distance $|n-c|$, which is unbounded. This proves the last assertion. We have reproduced the relevant argument of the supplied CATHOT C.5 \citep{CATHOTDraft2026}; no passage from a finite contraction to a uniform infinite contraction is used.
\end{proof}

\begin{corollary}[Proved: homotopical ordinary colimits of simplicial sets]\label{B:homotopical-colimits}
For a small nonempty connected $A$, total $\Winfty$-asphericity is equivalent to the following condition: for every $K:A\to\mathsf{sSet}$, the canonical augmentation from its simplicial-replacement homotopy colimit to its ordinary colimit is a weak equivalence. In particular, a total category's ordinary simplicial-set colimit functor preserves all objectwise weak equivalences.
\end{corollary}
\begin{proof}
The bisimplicial replacement has bidegree $(p,q)$ equal to
\[
 \coprod_{a_0\to\cdots\to a_p}K(a_0)_q.
\]
For a fixed $q$, this is the nerve of $\El_A(K_q)$, and its augmentation is the component map to $\colim_AK_q$. Theorem~\ref{C:classification} makes that augmentation a weak equivalence, component by component. The diagonal lemma for bisimplicial sets gives the asserted comparison for $K$. It follows by two-out-of-three that ordinary colimit preserves objectwise weak equivalences, since the simplicial replacement does: in each horizontal degree it is a coproduct of the given weak equivalences, and realization preserves these degreewise comparisons. Conversely apply the stated augmentation property to a discrete set-valued diagram $K=H$. Its source is $B\El_A(H)$ and its target is the discrete set of components, so Theorem~\ref{C:classification} gives totality. All comparison maps here are the canonical augmentations with the actual transition maps of $K$.
\end{proof}
